% GENERATED FILE. Edit canonical lessons, then run npm run generate:books.
% canonical-source-hash: fnv1a64:15b39ffe045b4b29
% canonical-lessons: GU-C296-embyulans, GU-C296-upadvu-lift, GU-C296-pato, GU-C296-tako, GU-C296-halat

\chapter{The Ambulance Arrives}
\label{ch:gu-the-ambulance-arrives}
\hlchaptermodality{296}

\begin{chapteropening}
\textbf{By the end of this chapter:} \emph{I can say how the injured were helped on the spot.}

Complete the last lesson of chapter 296: I can say how the injured were helped on the spot.
\end{chapteropening}

\section[\texorpdfstring{\gu{એમ્બ્યુલન્સ} (embyulans)}{એમ્બ્યુલન્સ (embyulans)}]{\gu{એમ્બ્યુલન્સ} (embyulans) — an ambulance}

\label{lesson:GU-C296-embyulans}



\begin{quote}
Before the new one: say the Gujarati for a fine, then the Gujarati for to crush.
\end{quote}

\subsection*{You'll want to know: \gu{એમ્બ્યુલન્સ}}
\textbf{\gu{એમ્બ્યુલન્સ}} — \emph{embyulans} — \textquotedblleft{}an ambulance\textquotedblright{}.

\begin{quote}
\gu{દસ} \gu{મિનિટમાં} \gu{એમ્બ્યુલન્સ} \gu{આવી}.
\end{quote}

\begin{quote}
\emph{das miniṭmā̃ embyulans āvī.}
\end{quote}

\begin{quote}
The ambulance came within ten minutes.
\end{quote}

A feminine noun.

\begin{cousinweb}
\textbf{\gu{એમ્બ્યુલન્સ}} (\emph{embyulans}) is borrowed from English \emph{ambulance}, and it is feminine: \textbf{\gu{એમ્બ્યુલન્સ} \gu{આવી}}, the ambulance came.
\end{cousinweb}

\begin{grammarlens}[title={What you've built}]
One more word for telling what happened.
\end{grammarlens}

\subsection*{Your turn}
\begin{itemize}
\raggedright
  \item \emph{Say it:} \emph{embyulans}
  \item \emph{Say it:} \emph{embyulans}, once more
  \item \emph{Say it:} \emph{das miniṭmā̃ embyulans āvī.} — The ambulance came within ten minutes.
  \item \emph{Recall:} say the Gujarati for a fine, then the Gujarati for to crush, then say \emph{embyulans} again
\end{itemize}

\begin{tcolorbox}[breakable,colback=teal!4,colframe=teal!35!black,title={Before you move on}]
What does \gu{એમ્બ્યુલન્સ} mean? (\textquotedblleft{}An ambulance\textquotedblright{}.) Say it once more.
\end{tcolorbox}
\section[\texorpdfstring{\gu{ઉપાડવું} (upāḍvũ)}{ઉપાડવું (upāḍvũ)}]{\gu{ઉપાડવું} (upāḍvũ) — to pick up; to carry off}

\label{lesson:GU-C296-upadvu-lift}



\begin{quote}
Before the new one: say the Gujarati for to crush, then the Gujarati for an ambulance.
\end{quote}

\subsection*{You'll want to know: \gu{ઉપાડવું}}
\textbf{\gu{ઉપાડવું}} — \emph{upāḍvũ} — \textquotedblleft{}to pick up; to carry off\textquotedblright{}.

\begin{quote}
\gu{અરુણે} \gu{અને} \gu{એક} \gu{સિપાઈએ} \gu{ઘાયલ} \gu{છોકરાને} \gu{ઉપાડ્યો}.
\end{quote}

\begin{quote}
\emph{Aruṇe ane ek sipāīe ghāyal chokrāne upāḍyo.}
\end{quote}

\begin{quote}
Arun and a constable picked up the injured boy.
\end{quote}

Past, with \gu{તેણે} (he or she …): \gu{તેણે} \gu{થેલી} \gu{ઉપાડી}, \emph{teṇe thelī upāḍī}.

\begin{cousinweb}
\textbf{\gu{ઉપાડવું}} (\emph{upāḍvũ}) goes back to Sanskrit \emph{útpātayati}, \textquotedblleft{}causes to fly up or rise\textquotedblright{}, and CDIAL lists the Gujarati verb, \textquotedblleft{}to lift up\textquotedblright{} (CDIAL 1810). It is lifting something to take it away: \textbf{\gu{ઘાયલ} \gu{છોકરાને} \gu{ઉપાડ્યો}}.
\end{cousinweb}

\begin{grammarlens}[title={What you've built}]
One more word for telling what happened.
\end{grammarlens}

\subsection*{Your turn}
\begin{itemize}
\raggedright
  \item \emph{Say it:} \emph{upāḍvũ}
  \item \emph{Say it:} \emph{upāḍvũ}, once more
  \item \emph{Say it:} \emph{Aruṇe ane ek sipāīe ghāyal chokrāne upāḍyo.} — Arun and a constable picked up the injured boy.
  \item \emph{Recall:} say the Gujarati for to crush, then the Gujarati for an ambulance, then say \emph{upāḍvũ} again
\end{itemize}

\begin{tcolorbox}[breakable,colback=teal!4,colframe=teal!35!black,title={Before you move on}]
What does \gu{ઉપાડવું} mean? (\textquotedblleft{}To pick up; to carry off\textquotedblright{}.) Say it once more.
\end{tcolorbox}
\section[\texorpdfstring{\gu{પાટો} (pāṭo)}{પાટો (pāṭo)}]{\gu{પાટો} (pāṭo) — a bandage}

\label{lesson:GU-C296-pato}



\begin{quote}
Before the new one: say the Gujarati for an ambulance, then the Gujarati for to pick up.
\end{quote}

\subsection*{You'll want to know: \gu{પાટો}}
\textbf{\gu{પાટો}} — \emph{pāṭo} — \textquotedblleft{}a bandage\textquotedblright{}.

\begin{quote}
\gu{નર્સે} \gu{તેના} \gu{હાથ} \gu{પર} \gu{પાટો} \gu{બાંધ્યો}.
\end{quote}

\begin{quote}
\emph{narse tenā hāth par pāṭo bā̃dhyo.}
\end{quote}

\begin{quote}
The nurse bandaged his arm.
\end{quote}

A masculine noun.

\begin{cousinweb}
\textbf{\gu{પાટો}} (\emph{pāṭo}) goes back to Sanskrit \emph{paṭṭa}, \textquotedblleft{}cloth, a bandage\textquotedblright{}, and CDIAL lists the Gujarati noun (CDIAL 7700). A bandage is tied on: \textbf{\gu{પાટો} \gu{બાંધ્યો}}.
\end{cousinweb}

\begin{grammarlens}[title={What you've built}]
One more word for telling what happened.
\end{grammarlens}

\subsection*{Your turn}
\begin{itemize}
\raggedright
  \item \emph{Say it:} \emph{pāṭo}
  \item \emph{Say it:} \emph{pāṭo}, once more
  \item \emph{Say it:} \emph{narse tenā hāth par pāṭo bā̃dhyo.} — The nurse bandaged his arm.
  \item \emph{Recall:} say the Gujarati for an ambulance, then the Gujarati for to pick up, then say \emph{pāṭo} again
\end{itemize}

\begin{tcolorbox}[breakable,colback=teal!4,colframe=teal!35!black,title={Before you move on}]
What does \gu{પાટો} mean? (\textquotedblleft{}A bandage\textquotedblright{}.) Say it once more.
\end{tcolorbox}
\section[\texorpdfstring{\gu{ટાંકો} (ṭā̃ko)}{ટાંકો (ṭā̃ko)}]{\gu{ટાંકો} (ṭā̃ko) — a stitch}

\label{lesson:GU-C296-tako}



\begin{quote}
Before the new one: say the Gujarati for to pick up, then the Gujarati for a bandage.
\end{quote}

\subsection*{You'll want to know: \gu{ટાંકો}}
\textbf{\gu{ટાંકો}} — \emph{ṭā̃ko} — \textquotedblleft{}a stitch\textquotedblright{}.

\begin{quote}
\gu{દાક્તરે} \gu{તેના} \gu{માથા} \gu{પર} \gu{ચાર} \gu{ટાંકા} \gu{લીધા}.
\end{quote}

\begin{quote}
\emph{dāktare tenā māthā par cār ṭā̃kā līdhā.}
\end{quote}

\begin{quote}
The doctor put four stitches in his head.
\end{quote}

A masculine noun.

\begin{cousinweb}
\textbf{\gu{ટાંકો}} (\emph{ṭā̃ko}) goes back to Sanskrit \emph{ṭaṅkáyati}, \textquotedblleft{}ties\textquotedblright{}, and CDIAL lists the Gujarati noun (CDIAL 5432). A doctor \textbf{\gu{ટાંકા} \gu{લે} \gu{છે}}, puts in stitches; a tailor takes them too.
\end{cousinweb}

\begin{grammarlens}[title={What you've built}]
One more word for telling what happened.
\end{grammarlens}

\subsection*{Your turn}
\begin{itemize}
\raggedright
  \item \emph{Say it:} \emph{ṭā̃ko}
  \item \emph{Say it:} \emph{ṭā̃ko}, once more
  \item \emph{Say it:} \emph{dāktare tenā māthā par cār ṭā̃kā līdhā.} — The doctor put four stitches in his head.
  \item \emph{Recall:} say the Gujarati for to pick up, then the Gujarati for a bandage, then say \emph{ṭā̃ko} again
\end{itemize}

\begin{tcolorbox}[breakable,colback=teal!4,colframe=teal!35!black,title={Before you move on}]
What does \gu{ટાંકો} mean? (\textquotedblleft{}A stitch\textquotedblright{}.) Say it once more.
\end{tcolorbox}
\section[\texorpdfstring{\gu{હાલત} (hālat)}{હાલત (hālat)}]{\gu{હાલત} (hālat) — a condition (of health or repair)}

\label{lesson:GU-C296-halat}



\begin{quote}
Before the new one: say the Gujarati for a bandage, then the Gujarati for a stitch.
\end{quote}

\subsection*{You'll want to know: \gu{હાલત}}
\textbf{\gu{હાલત}} — \emph{hālat} — \textquotedblleft{}a condition (of health or repair)\textquotedblright{}.

\begin{quote}
\gu{હવે} \gu{દર્દીની} \gu{હાલત} \gu{સારી} \gu{છે}.
\end{quote}

\begin{quote}
\emph{have dardīnī hālat sārī che.}
\end{quote}

\begin{quote}
The patient is in a good condition now.
\end{quote}

A feminine noun.

\begin{cousinweb}
\textbf{\gu{હાલત}} (\emph{hālat}) is usually traced, through Persian, to Arabic \emph{ḥāla}, \textquotedblleft{}a state, a condition\textquotedblright{}. It is feminine: \textbf{\gu{તેની} \gu{હાલત} \gu{ગંભીર} \gu{છે}}, his condition is serious.
\end{cousinweb}

\begin{grammarlens}[title={What you've built}]
One more word for telling what happened.
\end{grammarlens}

\subsection*{Your turn}
\begin{itemize}
\raggedright
  \item \emph{Say it:} \emph{hālat}
  \item \emph{Say it:} \emph{hālat}, once more
  \item \emph{Say it:} \emph{have dardīnī hālat sārī che.} — The patient is in a good condition now.
  \item \emph{Recall:} say the Gujarati for a bandage, then the Gujarati for a stitch, then say \emph{hālat} again
\end{itemize}

\begin{tcolorbox}[breakable,colback=teal!4,colframe=teal!35!black,title={Before you move on}]
What does \gu{હાલત} mean? (\textquotedblleft{}A condition (of health or repair)\textquotedblright{}.) Say it once more.
\end{tcolorbox}
